To recapitulate, the \emph{adaptivity} is only analyzed from the \emph{quantitative} perspective in previous works.
However, as we discussed before in the introduction, we consider it as a \emph{quantitative-reachability} property
that has two perspectives sub-properties, the  \emph{quantitative} and \emph{reachability}.
In the other word, previous works only analyzed it as a sub-property and lost the information from the other perspective.

\emph{Expressiveness.}
    This language model doesn't support programs with while loop whose number of iterations not known before loop is running.
    It only allows the while loops in the program to iterate a fixed number of times.
    This is because the expression in the loop guard has to be a constant number
    or an arithmetic expression
    that is evaluated into a constant $V_N$ before executing the loop body.
    However, the adaptive data analysis programs with loops of non-deterministic iteration times are very common.
    The syntax doesn't support the programs with user inputs either, which are also common in the data analysis area.

\emph{Accuracy.}
 The adaptivity definition and the estimated adaptivity are both inaccurate.
    \begin{itemize}
        \item 
        For the \emph{adaptivity} definition,
    %     There are two reasons for the inaccuracy of adaptivity definition.
    % The first reason comes from the trace in the operational semantics.
    % A trace only tracks the query requests, so it loses the execution history of the variables
    % that use the query results but not assigned by query requests.
    % These variables
    % are critical in analyzing the \emph{adaptivity}.
    % % \\
    % The second cause is the dependency graph definition.
    because the dependency graph relies on a specific memory, while-map, and a specific execution trace,
    the \emph{adaptivity} then is also defined w.r.t. one specific
    execution.
    In this sense, this adaptivity definition isn't the \emph{adaptivity} of the adaptive data analysis
    implemented by this program,
    but only its \emph{adaptivity} in a certain execution.
    \item
    The adaptivity bound computed from this framework is loose.
    Because the adaptivity estimation algorithm composes three standard static analysis techniques straightforwardly.
    The \textbf{graph generation} algorithm that uses the standard data-flow analysis technique
    over-approximates the data dependency relation in a large scale.
    And the standard path search algorithm also over-approximates the longest weighted path on this over-approximated dependency graph. 
    In this sense, the estimated adaptivity is a loose bound on the defined adaptivity.
 \end{itemize}

 \emph{Efficiency.}
 \begin{itemize}
    \item The operational semantics has four components in the configuration when evaluating the program. 
    % The updating operations for this quadruple configuration are inefficient, especially the update operation of the while map.   
    Updating this quadruple configuration (especially the while map) is inefficient.   
    \item
    The adaptivity estimation algorithm transforms the loop language into SSA form, in order to address the issue of re-assignment of different queries requesting results to the same variable.
    This transformation is inefficient and unnecessary.
    The re-assignment problem can be resolved efficiently and accurately through other static program analysis techniques, such as the
    variable reachable analysis, etc..
    \item The \textbf{variable estimation} and \textbf{graph generation} algorithms are inefficient.
    The \textbf{ag-loop} rule in Figure~\ref{fig:prework-static_alg1}
   and the \textbf{ad-loop} rule in Figure~\ref{fig:prework-static_alg2}
%    These two rules 
   unfold every iteration of the while loop.
    % and create new annotated variables for every iteration.
   This operation increased the algorithms' complexity exponentially. 
   %  \item For the same reason as above, their \textbf{Graph Generation} algorithm is low-efficient as well.
    \end{itemize}

Given the limitations in previous work, we develop our new framework
composed of three parts as well as in Figure~\ref{fig:adaptfun-overview}.
It is presented in the following three chapters.